% Solutions to the exercises for VI/03.
\begin{enumerate}[label=\textbf{E\arabic*.},leftmargin=*]
\item The restriction of \(f\) is zero iff \(f(a)=0\) for all \(a\in X\), exactly the condition \(f\in I(X)\).
\item The restrictions agree iff \((f-g)|_X=0\), which by E1 is equivalent to \(f-g\in I(X)\).
\item For infinite \(k\), a polynomial vanishing on all of \(k^n\) is zero, so \(I(\mathbb A_k^n)=(0)\).  Hence the quotient defining the coordinate ring is the polynomial ring itself.
\item One has \(I(\{a\})=(x_1-a_1,\dots,x_n-a_n)\).  Evaluation at \(a\) has this kernel and image \(k\), so the quotient is \(k\).
\item In the quotient, \([y]=(c-a[x])/b\).  Thus \([x]\) generates the ring.  Send \(t\mapsto[x]\); the inverse sends \([x]\mapsto t\), \([y]\mapsto(c-at)/b\).
\item Define \(\Phi:k[x,y]/(y-x^2)\to k[t]\) by \([x]\mapsto t\), \([y]\mapsto t^2\), and \(\Psi:k[t]\to k[x,y]/(y-x^2)\) by \(t\mapsto[x]\).  Since \([y]=[x]^2\), the compositions are identities.
\item The map \(k[t,t^{-1}]\to k[x,y]/(xy-1)\) with \(t\mapsto[x]\), \(t^{-1}\mapsto[y]\) is well-defined and surjective.  The reverse map sends \([x]\mapsto t\), \([y]\mapsto t^{-1}\); the relations are preserved, so the maps are inverse.
\item Every element of the quotient is represented by a polynomial in \(x_1,\dots,x_n\), hence by a polynomial in \(\bar x_i\).  A polynomial relation vanishes in the quotient iff its representative lies in the quotient kernel \(I(X)\).
\item If \([f]^m=0\), then \(f^m\in I(X)\).  Since \(I(X)\) is radical, \(f\in I(X)\), so \([f]=0\).
\item Evaluation respects sums, products, and constants.  It is onto because every \(c\in k\) is the image of the constant function \(c\).  The quotient by its kernel is \(k\), a field, so the kernel is maximal.
\item The natural map \(k[x]/I(X)\to k[x]/I(Y)\) is surjective and has kernel \(I(Y)/I(X)\).  Apply the first isomorphism theorem.
\item Since \((x-a)+(x-b)=(1)\), CRT gives \(k[x]/((x-a)(x-b))\cong k\times k\).  The idempotents are \(e_a=(x-b)/(a-b)\) and \(e_b=(x-a)/(b-a)\) modulo \((x-a)(x-b)\).
\item The ideals \((x-a_i)\) are pairwise comaximal and their intersection equals their product.  CRT gives \(k[X]\cong\prod_i k\cong k^r\).
\item The classes \(\bar x\) and \(\bar y\) are nonzero in \(k[x,y]/(xy)\), but \(\bar x\bar y=0\), so each is a zero divisor.
\item In \(B\), the class \(\bar x\neq0\) but \(\bar x^2=0\).  A classical coordinate ring is reduced, so \(B\) cannot be the coordinate ring of the one-point algebraic set.
\item Divide \(f\in\mathbb F_q[x]\) by \(x^q-x\): \(f=q_0(x)(x^q-x)+r\) with \(\deg r<q\).  If \(f\) vanishes at all \(q\) elements, then so does \(r\); a nonzero polynomial of degree below \(q\) cannot have \(q\) roots, so \(r=0\).  Thus the vanishing ideal is exactly \((x^q-x)\).
\item A finite set of \(rs\) rational points has coordinate ring equal to the algebra of all functions to \(k\), which is the product of \(rs\) copies of \(k\).
\item In \(k[X]\), \((\bar x+i\bar y)(\bar x-i\bar y)=\bar x^2+\bar y^2=1\).  Hence \(u^{-1}=\bar x-i\bar y\).
\end{enumerate}
